Fix \(g\).  We first identify the canonical linear program with the radius in \(\texttt{def:radius-objects}\).  Given a sign-reversing \(\mathbf K\in\mathfrak K_r(\mu^0,\nu)\), set
\(u_{g'xz}=|k_{g'x}(z)-k_x(z)|\).  These values satisfy every constraint of \(\mathsf P_g(\mu^0,\nu)\) with objective \(r\).  Conversely, primal feasibility implies
\[
 |k_{g'x}(z)-k_x(z)|\le u_{g'xz},
 \qquad
 \operatorname{TV}(k_{g'x},k_x)
 \le\frac12\sum_z u_{g'xz}\le r.
\]
Together with the stochastic-row, marginal-reproduction, and sign constraints, this says exactly that the row variables form a sign-reversing member of \(\mathfrak K_r(\mu^0,\nu)\).  Thus the feasible objective values are precisely radii in \([0,1]\) at which trial \(g\) can be reversed.

When the primal is feasible, its feasible set is compact: \(0\le r\le1\), every row coordinate lies in \([0,1]\), and \(0\le u_{g'xz}\le2r\le2\).  Hence its minimum is attained.  The equivalence just proved and upward nesting of \(\mathfrak K_r\) show that this minimum equals \(\delta_{\mathrm{sign},g}\).  Assumption \(\texttt{ass:positive-first-stage}\), through \(\texttt{lem:first-stage-denominator}\), ensures that this numerator sign is also the sign of the PSACE and that no zero denominator is being classified.

Apply the finite-dimensional linear-program alternative in \(\texttt{lem:finite-linear-program-certificates}\) to the standard form from \(\texttt{def:sign-radius-dual}\).  For primal-feasible \(w\) and dual-feasible \((y,\lambda)\), weak duality is visible directly:
\[
 e(\nu)^{\top}y-f^{\top}\lambda
 \le e(\nu)^{\top}y-\lambda^{\top}F_gw
 =(E_\mu^{\top}y-F_g^{\top}\lambda)^{\top}w
 \le c^{\top}w.
\]
Strong duality gives equality at optimizers.  After quotienting the declared lineality, the dual is pointed.  A finite optimum of a pointed polyhedron is attained at a basic feasible point; there are only finitely many bases.  It follows that
\[
 \delta_{\mathrm{sign},g}
 =\min\mathsf P_g(\mu^0,\nu)
 =\max_{(y,\lambda)\in\mathcal B_g(\mu^0)}
   \Phi_g(\mu^0,\nu;y,\lambda)\in[0,1].
\]
The primal optimizer itself supplies the advertised adversarial kernel.  The maximizing dual cut has value \(r_*\); weak duality shows that every sign-reversing kernel has radius at least \(r_*\), so it is the advertised lower-radius certificate.

If the primal is infeasible, the first equivalence says that no sign-reversing compatible system exists for any \(r\le1\), including \(r=1\).  Hence the defining set for \(\delta_{\mathrm{sign},g}\) is empty and the stipulated convention gives \(+\infty\).  The infeasible alternative in \(\texttt{lem:finite-linear-program-certificates}\) returns a Farkas multiplier whose checked linear combination makes the right side strictly negative while the left side is nonnegative; it therefore certifies the same impossibility.

If the marginals are rational, every coefficient in \(E_\mu,e(\nu),F_g,f,c\) is rational.  The cited linear-program alternative permits a primal optimizer and dual basic optimizer obtained by solving nonsingular rational subsystems, hence rational, and in the infeasible case a rational extreme Farkas ray.  Finally, because \(\mathcal G\) is finite,
\(\min_g\delta_{\mathrm{sign},g}=+\infty\) if and only if each term is \(+\infty\), which by the preceding equivalence is exactly the case in which every trial has its infeasibility certificate.